\documentclass[10pt,a4paper]{amsart}
\usepackage{amsmath,amssymb,amsthm,graphicx}
% Retained preamble for the Stacks Project chapter snapshot in this directory.
% The chapter's own preamble.tex is a copy of the shared Stacks Project preamble; the block
% below is sliced MECHANICALLY, line for line, from it. Nothing was added, removed or
% redefined; the wrapper class, the xr-hyper/hyperref wiring and the externaldocument list
% are omitted because the assembler supplies the document class.
%
% lines taken from preamble.tex: 3, 5, 9, 11, 12, 156, 157, 158, ... (51 lines total)
%
% package coverage restored from the chapter's own preamble (2026-10-13): xy, xy-2cell

\usepackage[all]{xy}
\xyoption{2cell}
\UseAllTwocells

\newenvironment{reference}{\comment}{\endcomment}
\newenvironment{slogan}{\comment}{\endcomment}
\newenvironment{history}{\comment}{\endcomment}
\theoremstyle{plain}
\newtheorem{theorem}[subsection]{Theorem}
\newtheorem{proposition}[subsection]{Proposition}
\newtheorem{lemma}[subsection]{Lemma}
\theoremstyle{definition}
\newtheorem{definition}[subsection]{Definition}
\newtheorem{example}[subsection]{Example}
\newtheorem{exercise}[subsection]{Exercise}
\newtheorem{situation}[subsection]{Situation}
\theoremstyle{remark}
\newtheorem{remark}[subsection]{Remark}
\newtheorem{remarks}[subsection]{Remarks}
\numberwithin{equation}{subsection}
\def\lim{\mathop{\mathrm{lim}}\nolimits}
\def\colim{\mathop{\mathrm{colim}}\nolimits}
\def\Spec{\mathop{\mathrm{Spec}}}
\def\Hom{\mathop{\mathrm{Hom}}\nolimits}
\def\Ext{\mathop{\mathrm{Ext}}\nolimits}
\def\SheafHom{\mathop{\mathcal{H}\!\mathit{om}}\nolimits}
\def\SheafExt{\mathop{\mathcal{E}\!\mathit{xt}}\nolimits}
\def\Sch{\mathit{Sch}}
\def\Mor{\mathop{\mathrm{Mor}}\nolimits}
\def\Ob{\mathop{\mathrm{Ob}}\nolimits}
\def\Sh{\mathop{\mathit{Sh}}\nolimits}
\def\NL{\mathop{N\!L}\nolimits}
\def\CH{\mathop{\mathrm{CH}}\nolimits}
\def\proetale{{pro\text{-}\acute{e}tale}}
\def\etale{{\acute{e}tale}}
\def\QCoh{\mathit{QCoh}}
\def\Ker{\mathop{\mathrm{Ker}}}
\def\Im{\mathop{\mathrm{Im}}}
\def\Coker{\mathop{\mathrm{Coker}}}
\def\Coim{\mathop{\mathrm{Coim}}}
\DeclareMathSymbol{\boxtimes}{\mathbin}{AMSa}{"02}
\def\QCohstack{\mathcal{QC}\!\mathit{oh}}
\def\Cohstack{\mathcal{C}\!\mathit{oh}}
\def\Spacesstack{\mathcal{S}\!\mathit{paces}}
\def\Quotfunctor{\mathrm{Quot}}
\def\Hilbfunctor{\mathrm{Hilb}}
\def\Curvesstack{\mathcal{C}\!\mathit{urves}}
\def\Polarizedstack{\mathcal{P}\!\mathit{olarized}}
\def\Complexesstack{\mathcal{C}\!\mathit{omplexes}}
\def\Pic{\mathop{\mathrm{Pic}}\nolimits}
\def\Picardstack{\mathcal{P}\!\mathit{ic}}
\def\Picardfunctor{\mathrm{Pic}}
\def\Deformationcategory{\mathcal{D}\!\mathit{ef}}
\usepackage{cleveref}
\crefname{originallemma}{original lemma}{original lemmas}
\Crefname{originallemma}{Original lemma}{Original lemmas}
\makeatletter
\newlabel{algebra-lemma-differential-seq}{{00RU}{source}{}{sourcepointer-1}{}}
\newlabel{algebra-lemma-differential-seq@cref}{{[originallemma][0][]00RU}{[1][0][]source}{}{}{}}
\makeatother
\crefname{originallemma}{original lemma}{original lemmas}
\Crefname{originallemma}{Original lemma}{Original lemmas}
\makeatletter
\newlabel{algebra-lemma-differential-surjective}{{00RR}{source}{}{sourcepointer-2}{}}
\newlabel{algebra-lemma-differential-surjective@cref}{{[originallemma][0][]00RR}{[1][0][]source}{}{}{}}
\makeatother
\crefname{originallemma}{original lemma}{original lemmas}
\Crefname{originallemma}{Original lemma}{Original lemmas}
\makeatletter
\newlabel{lemma-compute-cotangent-complex}{{08PU}{source}{}{sourcepointer-3}{}}
\newlabel{lemma-compute-cotangent-complex@cref}{{[originallemma][0][]08PU}{[1][0][]source}{}{}{}}
\makeatother
\crefname{originalsection}{original section}{original sections}
\Crefname{originalsection}{Original section}{Original sections}
\makeatletter
\newlabel{sites-cohomology-example-category-to-point}{{08PF}{source}{}{sourcepointer-4}{}}
\newlabel{sites-cohomology-example-category-to-point@cref}{{[originalsection][0][]08PF}{[1][0][]source}{}{}{}}
\makeatother
\title{The zeroth cohomology of the cotangent complex is the module of differentials}
\author{The Stacks Project Authors}
\date{stacks-project:cotangent.tex:local-snapshot-20261012; GNU FDL 1.2 or later}
\begin{document}
\maketitle
\setcounter{subsection}{1}
\noindent\textit{Editorial note.} All delivered mathematical text is the author's own, copied from the pinned Stacks Project chapter file cotangent.tex; no statement or proof step was rewritten, reordered or removed.\par\medskip
\noindent\textit{Editorial note.} The statement is one sentence; the proof is the author's explicit two-part calculation (surjectivity and injectivity) inside the category C\_\{B/A\}.\par\medskip
\noindent\textit{Editorial note.} The four references of the counted spans are recorded as non-load-bearing navigation pointers; the argument that consumes them stays inside the counted proof span.\par\medskip
\noindent\textit{Editorial note.} Preamble: the retained preamble of this chapter (sources/stacks-cotangent/retained-preamble.tex, lines 1-63), containing the Xy-pic setup, the comment environments, the theorem declarations and the macros including \textbackslash{}NL.\par\medskip
\noindent\textit{Allowed standard prerequisites (editorial).}\begin{itemize}
\item The identification of the cotangent complex as Lpi\_! of the differentials of the standard resolution. Lemma lemma-compute-cotangent-complex, used at proof lines 588-590. Turns the claim into a statement about pi\_! of a sheaf on C\_\{B/A\}.
\item Surjectivity of the differential onto the module of differentials. Algebra, Lemma algebra-lemma-differential-surjective, used at proof lines 598-600. Gives surjectivity of the canonical map.
\item The exact sequence for the differentials of a quotient ring. Algebra, Lemma algebra-lemma-differential-seq, used at proof lines 604-609. Identifies the kernel with the image of I/I\textasciicircum{}2.
\item The colimit description of pi\_! for the punctual topos. Cohomology on Sites, Example sites-cohomology-example-category-to-point, used at proof lines 614-616. Converts the two agreeing images into the injectivity statement.
\end{itemize}\medskip
\ifdefined\hypertarget\hypertarget{sourcepointer-1}{}\fi
\noindent\textit{Source reference.} Original lemma 00RU: focused-seeds/AGP/sources/stacks-snapshot-20261011/tags.txt, source lines 1016--1016. Navigation pointer into the Stacks Project; the standard fact it supplies is applied inside the counted proof span and the argument resting on it stays there.\par\medskip
\ifdefined\hypertarget\hypertarget{sourcepointer-2}{}\fi
\noindent\textit{Source reference.} Original lemma 00RR: focused-seeds/AGP/sources/stacks-snapshot-20261011/tags.txt, source lines 1013--1013. Navigation pointer into the Stacks Project; the standard fact it supplies is applied inside the counted proof span and the argument resting on it stays there.\par\medskip
\ifdefined\hypertarget\hypertarget{sourcepointer-3}{}\fi
\noindent\textit{Source reference.} Original lemma 08PU: focused-seeds/AGP/sources/stacks-snapshot-20261011/tags.txt, source lines 11037--11037. Navigation pointer into the Stacks Project; the standard fact it supplies is applied inside the counted proof span and the argument resting on it stays there.\par\medskip
\ifdefined\hypertarget\hypertarget{sourcepointer-4}{}\fi
\noindent\textit{Source reference.} Original section 08PF: focused-seeds/AGP/sources/stacks-snapshot-20261011/tags.txt, source lines 11023--11023. Navigation pointer into the Stacks Project; the standard fact it supplies is applied inside the counted proof span and the argument resting on it stays there.\par\medskip
\begin{lemma}
\label{lemma-identify-H0}
If $A \to B$ is a ring map, then $H^0(L_{B/A}) = \Omega_{B/A}$.
\end{lemma}

\begin{proof}
We will prove this by a direct calculation.
We will use the identification of Lemma \ref{lemma-compute-cotangent-complex}.
There is clearly a map from $\Omega_{\mathcal{O}/A} \otimes \underline{B}$
to the constant sheaf with value $\Omega_{B/A}$. Thus this map induces
a map
$$
H^0(L_{B/A}) = H^0(L\pi_!(\Omega_{\mathcal{O}/A} \otimes \underline{B}))
= \pi_!(\Omega_{\mathcal{O}/A} \otimes \underline{B}) \to \Omega_{B/A}
$$
By choosing an object $P \to B$ of $\mathcal{C}_{B/A}$ with $P \to B$
surjective we see that this map is surjective (by
Algebra, Lemma \ref{algebra-lemma-differential-surjective}).
To show that it is injective, suppose that $P \to B$ is an object
of $\mathcal{C}_{B/A}$ and that $\xi \in \Omega_{P/A} \otimes_P B$
is an element which maps to zero in $\Omega_{B/A}$.
We first choose factorization $P \to P' \to B$ such that $P' \to B$
is surjective and $P'$ is a polynomial algebra over $A$.
We may replace $P$ by $P'$. If $B = P/I$, then the kernel
$\Omega_{P/A} \otimes_P B \to \Omega_{B/A}$ is the image of
$I/I^2$ (Algebra, Lemma \ref{algebra-lemma-differential-seq}).
Say $\xi$ is the image of $f \in I$.
Then we consider the two maps $a, b : P' = P[x] \to P$, the first of which
maps $x$ to $0$ and the second of which maps $x$ to $f$ (in both
cases $P[x] \to B$ maps $x$ to zero). We see that $\xi$ and $0$
are the image of $\text{d}x \otimes 1$ in $\Omega_{P'/A} \otimes_{P'} B$.
Thus $\xi$ and $0$ have the same image in the colimit (see
Cohomology on Sites, Example \ref{sites-cohomology-example-category-to-point})
$\pi_!(\Omega_{\mathcal{O}/A} \otimes \underline{B})$ as desired.
\end{proof}

\end{document}
